\section{Background}
\label{sec:background}

\paragraph{REST APIs and OpenAPI.} A REST API offers \emph{operations}: an
HTTP method on a path template, such as \texttt{POST /owners/\{ownerId\}/pets}.
An \emph{OpenAPI specification}~\cite{openapi} is a JSON or YAML document that
lists every operation with its parameters (in the path, query, header or
cookie), its request body and its responses. Each parameter and body has a
\emph{schema}: a type plus optional limits such as \texttt{enum},
\texttt{minimum}, \texttt{maxLength}, a regular expression in
\texttt{pattern}, or a \texttt{format} such as \texttt{date}.

\paragraph{References and composition.} A schema is rarely written in place.
It points to a shared definition with \texttt{\$ref}, for example
\texttt{\{"\$ref": "\#/components/schemas/Pet"\}}. Parameters, request bodies
and responses can be references too. A schema can also be built from parts:
\texttt{allOf} means ``all of these at once'', while \texttt{oneOf} and
\texttt{anyOf} offer alternatives. A \emph{recursive schema} reaches itself
through references, such as a tree node whose children are tree nodes.

\paragraph{Black-box REST API testing.} A black-box tester sees only the
specification and the HTTP responses, not the source code. It sends requests
and reads the status code. A 2XX status means the API accepted the request.
A 4XX status means the API rejected it as a client error. A 5XX status means
the server failed. Testers usually run two kinds of requests: \emph{nominal}
requests that try to follow the specification, and \emph{error} requests that
break it on purpose~\cite{resttestgen-stvr}.

\paragraph{Operation dependencies.} Many operations need a resource that
another operation creates. \texttt{GET /owners/\{ownerId\}} returns 404 unless
\texttt{ownerId} names an owner that exists. A tester must therefore create
resources first and reuse the identifiers (\emph{ids}) that come back.
RESTler calls this producer--consumer inference~\cite{restler}.

\paragraph{REST League.} REST League 2027~\cite{restleague} is a competition
for black-box REST API testing tools. Each tool gets one hour per API inside
RESTgym~\cite{restgym}. It is scored on (i)~unique 5XX responses, where two
5XX responses count once when their error messages are not sufficiently
distinct; (ii)~\emph{operation coverage}, the operations that returned at
least one 2XX; (iii)~code coverage; and (iv)~how fast the tool reaches these
values. Five of its ten APIs are known in advance; pet-clinic and
gestao-hospital, the two APIs we test live in Section~\ref{sec:eval}, are
among them.

\paragraph{Our team tool.} We build a REST League tool in three parts with a
fixed interface between them. The \emph{spec component} reads the
specification and produces requests. The \emph{engine} sends the requests,
keeps a \emph{resource pool} of ids seen in responses, and decides what to
try next. The \emph{packaging} runs the tool in RESTgym. This report is about
the spec component, which we call \tool. The engine calls \tool through one
function, \texttt{generate\_request(op, pool, mode)}, and \tool asks the pool
for real ids through \texttt{pool.lookup(name, schema, op)}.
